%==============================================================================
% introduction.tex — Introduction for PP-HFM
% Self-contained \section{}; \input this into main.tex (never compile alone).
%
% Structured after the house introduction template: domain motivation, the
% task, what makes it hard, how approaches developed, what each leaves open,
% the remedies that follow, "Accordingly, this work...", an enumerated
% contribution list with explicit non-claims, and the organisation paragraph.
%
% POSITIONING RULE (docs/roadmap.md): user-level DP, parameter-efficient DP
% fine-tuning and DP quantile release are PRIOR ART. The contributions are
% measurements and auditing practice. Do not write "we propose patient-level
% DP".
%
% NOTE: the contribution list covers only completed work. The records-per-
% individual scaling law completed 2026-08-16 (24/24 configurations) and is
% now the fifth contribution; see Sec. Results, "Records per individual".
% It is stated in AUPRC only -- the AUROC trend is monotone in the same
% direction but every interval covers zero at three seeds, so it is reported
% as directionally consistent and not as a resolved effect.
%==============================================================================

\section{Introduction}
\label{sec:intro}

Internet-of-Things deployments increasingly derive their value from combining
sensing streams of different kinds. A hospital theatre records continuous
physiological waveforms alongside the structured clinical history of the
patient on the table; a wearable installation records inertial motion alongside
egocentric video. In each case the streams are complementary rather than
redundant, and a model with access to both can resolve states that neither
resolves alone. The streams are also produced where the sensors are, which is
seldom where computation is convenient, and this tension between where data
arises and where models are trained is the setting the present work addresses.

Transmitting raw multimodal streams to a central service is the obvious
resolution and is frequently unavailable. The volume is substantial and the
uplink is the scarce resource, so centralisation inflates both latency and
energy in precisely the deployments that can least afford either. More
decisively, in the settings where multimodal sensing is most valuable the data
cannot lawfully move at all: clinical records are governed by consent and by
data-protection regimes that prohibit pooling across institutions, so a model
requiring centralisation cannot be deployed rather than merely being expensive.
Federated learning \cite{ref1} answers this by communicating model updates
instead of records, and hierarchical variants \cite{ref2, ref3} interpose an
edge tier so that the wide-area link is exercised once per several local
rounds.

Moving the model rather than the data does not by itself protect the data. The
updates a client transmits are functions of its records, and membership
inference against released models is well established \cite{ref14}. A formal
guarantee is therefore required, and differentially private stochastic gradient
descent \cite{ref6} supplies one by bounding the influence of a training record
and perturbing the aggregate. But the guarantee such a system reports is stated
over the unit it clips, and in a sensing deployment that unit is a
fixed-length window rather than a person. Where an individual contributes many
windows --- as every patient does, and as every participant in a continuous
recording does --- a record-level budget answers whether a particular window
was used in training, while the question a clinical deployment must answer is
whether a particular \emph{person} was. The two can be related by group
privacy \cite{ref8}, but at the group sizes these datasets exhibit the
conversion is prohibitive, turning a nominal budget of unity into one that
supports no meaningful claim.

Existing systems have developed along two lines that do not meet. Hierarchical
designs establish the communication and energy case convincingly: HED-FL
\cite{ref3} extends deployment lifetime by avoiding long-range transmission,
HCEF \cite{ref4} adapts computation and compression to device heterogeneity,
and assignment-oriented work \cite{ref5} reduces divergence under non-identical
distributions. None carries a formal privacy guarantee. Private designs supply
the guarantee: FL2DP \cite{ref9} couples noise with shuffling, and RoPPFL
\cite{ref10} pairs local privacy with robust aggregation at the edge. Neither
is hierarchical in the sense that matters here, and neither examines how an
intermediate tier alters the frequency with which each client's budget is
spent. Across both lines the local objective is treated as the aggregation of
generic parameter vectors, so neither multiple modalities per record nor the
grouping of records under an individual plays any part.

Three difficulties follow, and they bear on deployability as much as on
privacy. The unit at which a guarantee is stated is rarely examined, so its
cost on genuinely paired multimodal data is unquantified. The validity
condition of parameter-efficient private training --- that the corpus used to
pretrain the frozen representation be disjoint from the protected population
\cite{ref12, ref13} --- is frequently left implicit and, where honoured,
unpriced. And guarantees are typically reported from the parameters of the
mechanism rather than audited by attack; where an audit is performed it is
conducted at the granularity of the record, which is precisely the quantity a
record-level guarantee already protects.

Two remedies follow directly. The unit can be fixed at the point of clipping
rather than recovered afterwards by a group-privacy conversion: if the
gradients of an individual's records are accumulated before the bound is
applied, that individual's influence on the released model is bounded
irrespective of how many records they contributed, and accounting performed at
the same granularity yields a budget that is a statement about people by
construction. And the resulting guarantee can be audited by mounting membership
inference at the level of the individual, provided the attack used is strong
enough that a negative finding is informative --- which requires demonstrating
on a configuration known to memorise that the attack can find leakage where
leakage exists.

Accordingly, this work presents PP-HFM, a three-tier federated framework for
multimodal sensing in which the privacy unit is the individual by construction.
Records are grouped by individual before they are bounded; the noise scale is
calibrated to satisfy the budget for every client simultaneously rather than
for a representative one, so that the interposition of the edge tier is
reflected in the reported figure; and the cost of releasing a conformal
calibration quantile is composed into that figure rather than omitted from it.
Representation learning is separated from private adaptation, with encoders
pretrained on a population disjoint from every client and frozen, so that only
a compact fusion head is privatised. The framework is evaluated on two corpora
whose modalities are associated by the recording apparatus rather than by a
construction rule, under a protocol that measures communication, latency,
reliability and energy on a simulated topology and audits the privacy guarantee
by attack.

The contributions are fivefold. First, the framework enforces an
individual-level guarantee within a hierarchical multimodal system and reports
a composed budget covering both training and calibration; the accounting
practices this requires --- calibrating the noise against every client rather
than a representative one, and charging the calibration release against the
same budget --- are specified in enough detail to be checked from outside,
because each corresponds to a failure that is invisible in a training log.
Second, the cost of the choice of unit is measured on genuinely paired clinical
multimodal data, together with the cost of the disjointness condition that
makes the guarantee valid end-to-end, so that both are priced rather than
assumed negligible. Third, the guarantee is audited by patient-level membership
inference under both a threshold attack and a likelihood-ratio attack
\cite{ref17}, each reported alongside a positive control on a configuration
known to memorise; the audit returns a null on the private arms, and the
control is what makes that null attributable to the defence rather than to an
inadequate adversary. Fourth, the system-level behaviour is validated on two
deployments differing in topology and in payload, with every network quantity
measured on rounds in which the payload traversed the simulated topology rather
than composed analytically. The same protocol establishes that the guarantees
are affordable in real time: the conformal set construction that carries the
coverage guarantee occupies $0.68\,\%$ of the inference path, and the private
mechanism occupies none of it, since its noise is applied to gradients during
training and adds no operation to serving. The objection that a privacy
apparatus is too heavy for a latency-bound deployment is therefore answered by
measurement rather than by argument. Fifth, the price of the individual-level guarantee
is shown not to be a fixed quantity but a function of how much each individual
contributes: with the per-step budget and the number of units per step held
constant, the penalty is indistinguishable from zero when an individual
contributes a single record and grows monotonically as that number rises,
saturating by eight records. This design carries its own control, because at one
record per individual the two units coincide by construction and the two arms
must agree; that they do is what licenses reading the divergence at larger
contributions as a property of the unit rather than of the apparatus.

What is not claimed should be stated as plainly. User-level differential
privacy \cite{ref11}, parameter-efficient differentially private fine-tuning
\cite{ref12, ref13} and differentially private quantile release \cite{ref24}
are established prior art, and are used here rather than proposed. The
contributions above are measurements of these mechanisms in a setting where
they have not been measured, and the auditing practice under which those
measurements are reported.

The remainder of the paper is organised as follows. Section~\ref{sec:related}
reviews related work on hierarchical and private federated learning, on the
unit at which a guarantee is stated, and on membership inference as an
auditing instrument; Section~\ref{sec:method} presents the proposed framework;
Section~\ref{sec:experiments} specifies the corpora, the implementation, the
evaluation measures and the experimental protocol; Section~\ref{sec:results}
reports the results; Section~\ref{sec:discussion} discusses their
interpretation and the limitations of the study; and
Section~\ref{sec:conclusion} concludes.
